\documentclass[11pt,a4paper]{article}
\usepackage{fontspec,amsmath,amssymb,amsthm,microtype}
\setmainfont{lmroman10-regular.otf}[BoldFont=lmroman10-bold.otf,ItalicFont=lmroman10-italic.otf,BoldItalicFont=lmroman10-bolditalic.otf]
\usepackage{xcolor}
\usepackage[most]{tcolorbox}
\usepackage{mathrsfs}
\providecommand{\lt}{<}
\providecommand{\gt}{>}
\setmonofont{DejaVuSansMono.ttf}[BoldFont=DejaVuSansMono-Bold.ttf,ItalicFont=DejaVuSansMono-Oblique.ttf,BoldItalicFont=DejaVuSansMono-BoldOblique.ttf,Scale=0.86]
\newfontfamily\leanfallback{STIXTwoMath-Regular.otf}[Scale=0.86]
\definecolor{leanrule}{gray}{0.62}
\definecolor{leanback}{gray}{0.962}
\definecolor{leantitle}{gray}{0.22}
\definecolor{leanmuted}{gray}{0.38}
\newtcolorbox{leanbox}[1]{enhanced,breakable,colback=leanback,colframe=leanrule,
  boxrule=0.4pt,arc=0.6mm,left=2.4mm,right=2.4mm,top=1.8mm,bottom=1.8mm,
  before skip=0.8em,after skip=1em,fonttitle=\small\bfseries,coltitle=leantitle,
  colbacktitle=leanback,titlerule=0.3pt,title={#1}}
\newlength{\leancharwidth}
\newenvironment{leancode}{\par\smallskip\begingroup\ttfamily\footnotesize\raggedright
  \settowidth{\leancharwidth}{M}\setlength{\parskip}{0pt}\setlength{\parindent}{0pt}
  \frenchspacing}{\par\endgroup\smallskip}
\newcommand{\leanline}[2]{\hangindent=\dimexpr#1\leancharwidth+4\leancharwidth\relax
  \hangafter=1\noindent\hspace*{#1\leancharwidth}#2\par}
\newcommand{\leansource}[1]{\par{\footnotesize\color{leanmuted}#1}\par}
\newcommand{\leanlabel}[1]{\par\smallskip\noindent{\small\bfseries #1}\par}
\usepackage[margin=25mm]{geometry}
\usepackage[unicode,hidelinks]{hyperref}
\hypersetup{pdfauthor={LangevinLean Project}}
\setlength{\emergencystretch}{3em}
\setlength{\parskip}{0.4em}
\setlength{\parindent}{0pt}
\allowdisplaybreaks[2]
\raggedbottom
\title{\texorpdfstring{Bounded fields and positive global strong monotonicity}{Bounded fields and positive global strong monotonicity}}
\author{LangevinLean Project}
\date{}
\begin{document}\maketitle
Draft mathematical exposition; these notes may contain mistakes.


A globally bounded vector field on a nontrivial real Hilbert space cannot have a positive global strong-monotonicity constant. The Moreau gradient of a globally Lipschitz convex function is bounded. Hence the conjunction of full-potential value Lipschitzness and a positive global proximal-residual deficit is inconsistent. The conclusion concerns that literal conjunction of assumptions; it does not apply to a decomposition in which only the nonsmooth summand is Lipschitz.


\section{\texorpdfstring{Incompatibility of boundedness and positive global monotonicity}{Incompatibility of boundedness and positive global monotonicity}}


Assume \(\|R(x)\|\le K\) and \(c\|x-y\|^2\le\langle x-y,R(x)-R(y)\rangle\) with \(c>0\). Cauchy–Schwarz and the triangle inequality give \(c\|x-y\|^2\le2K\|x-y\|\). A nonzero direction can be scaled to have length greater than \(2K/c\), contradicting this bound.



In the alternative notation \(\|b(x)\|\le K\) and \(\langle b(x)-b(y),x-y\rangle\ge q\|x-y\|^2\) with \(q>0\) imply \(q\|x-y\|\le2K\) for every distinct pair. Explicitly choose separation \(r>2K/c\). Strong monotonicity bounds the inner product below by \(cr^2\), while boundedness bounds it above by \(2Kr\). Division by \(r\) gives \(cr\le2K\). A nontrivial space has arbitrarily large such separations \(r\), so the assumptions cannot both hold. The nontrivial-space qualification is essential.

\begin{leanbox}{Lean: \texttt{bounded\_\allowbreak{}field\_\allowbreak{}not\_\allowbreak{}stronglyMonotone}}
\leansource{\texttt{bounded\_\allowbreak{}field\_\allowbreak{}not\_\allowbreak{}stronglyMonotone} \quad ProximalLangevin/\allowbreak{}PIPULAConsistency.lean, lines 23–42}
\begin{leancode}
\leanline{0}{theorem\ bounded\_field\_not\_stronglyMonotone}
\leanline{4}{\{E\ :\ Type*\}\ [NormedAddCommGroup\ E]\ [InnerProductSpace\ ℝ\ E]\ [Nontrivial\ E]}
\leanline{4}{(R\ :\ E\ →\ E)\ (c\ K\ :\ ℝ)\ (hc\ :\ 0\ <\ c)\ (hK\ :\ 0\ ≤\ K)}
\leanline{4}{(hbound\ :\ ∀\ x,\ ‖R\ x‖\ ≤\ K)}
\leanline{4}{(hmono\ :\ ∀\ x\ y,\ c\ *\ ‖x\ -\ y‖\ \textasciicircum{}\ 2\ ≤\ {\leanfallback ⟪}x\ -\ y,\ R\ x\ -\ R\ y{\leanfallback ⟫})\ :\ False}
\end{leancode}
\leanlabel{How the Lean proof follows the argument}
Formal correspondence: bounded\_field\_not\_stronglyMonotone : A globally bounded field cannot have a positive strong-monotonicity constant on a nontrivial real Hilbert space.



\texttt{bounded\_\allowbreak{}field\_\allowbreak{}not\_\allowbreak{}strongly\allowbreak{}Monotone}: the intermediate facts include \texttt{r}, \texttt{hr}, \texttt{h\allowbreak{}R}, \texttt{hi}, \texttt{hrc}, in source order; the source steps apply or rewrite \texttt{sub\_\allowbreak{}zero}. The remaining scalar rearrangements are checked by \texttt{field\_\allowbreak{}simp}, \texttt{linarith}, \texttt{nlinarith}.
\end{leanbox}

\section{\texorpdfstring{The obstruction for a globally bounded proximal residual}{The obstruction for a globally bounded proximal residual}}


Consider the reading of the source reduction in which \(g_1=0\) also applies the value-Lipschitz assumption to \(U=g_2\). Then \(U\) is globally \(K\)-Lipschitz, so its Moreau gradient \(R_\gamma=\gamma^{-1}(I-P_\gamma)\) satisfies \(\|R_\gamma(x)\|\le K\). A positive proximal deficit gives

\[
\frac{\mu}{\gamma}\|x-y\|^2
 \le\langle x-y,R_\gamma(x)-R_\gamma(y)\rangle
 \le2K\|x-y\|.
\]

The preceding obstruction contradicts these simultaneous global conditions on a nontrivial space. If the Lipschitz assumption instead concerns the nonsmooth term in an independent decomposition \(U=g_1+g_2\), it does not make the full \(U\) Lipschitz and the contradiction does not follow.



The residual \(R_\gamma=(I-P_\gamma)/\gamma\) is bounded in the first reading, while the deficit makes it strongly monotone with coefficient \(\mu/\gamma\). Under \(\mu>0\) these are incompatible. This is a conditional consistency result about the specified assumptions, not a general refutation of the source corollary. The separate full-potential PIPULA theorem assumes strong convexity and a lower bound without global value Lipschitzness.

\begin{leanbox}{Lean: \texttt{lipschitz\_\allowbreak{}full\_\allowbreak{}potential\_\allowbreak{}not\_\allowbreak{}proximal\_\allowbreak{}deficit}}
\leansource{\texttt{lipschitz\_\allowbreak{}full\_\allowbreak{}potential\_\allowbreak{}not\_\allowbreak{}proximal\_\allowbreak{}deficit} \quad ProximalLangevin/\allowbreak{}PIPULAConsistency.lean, lines 47–57}
\begin{leancode}
\leanline{0}{theorem\ lipschitz\_full\_potential\_not\_proximal\_deficit}
\leanline{4}{\{E\ :\ Type*\}\ [NormedAddCommGroup\ E]\ [InnerProductSpace\ ℝ\ E]}
\leanline{4}{[FiniteDimensional\ ℝ\ E]\ [Nontrivial\ E]}
\leanline{4}{\{U\ :\ E\ →\ ℝ\}\ (hU\ :\ ConvexOn\ ℝ\ univ\ U)\ (hb\ :\ ∃\ b,\ ∀\ x,\ b\ ≤\ U\ x)}
\leanline{4}{\{K\ :\ ℝ≥0\}\ (hK\ :\ LipschitzWith\ K\ U)\ (μ\ γ\ :\ ℝ)\ (hμ\ :\ 0\ <\ μ)\ (hγ\ :\ 0\ <\ γ)}
\leanline{4}{(hB\ :\ ∀\ x\ y,\ {\leanfallback ⟪}x\ -\ y,\ prox\ U\ γ\ x\ -\ prox\ U\ γ\ y{\leanfallback ⟫}\ ≤}
\leanline{6}{(1\ -\ μ)\ *\ ‖x\ -\ y‖\ \textasciicircum{}\ 2)\ :\ False}
\end{leancode}
\leanlabel{How the Lean proof follows the argument}
Formal correspondence: lipschitz\_full\_potential\_not\_proximal\_deficit : Under a literal full-potential globally Lipschitz reading, B1 cannot hold with positive deficit in positive dimension. The separate PIPULA theorem therefore does not assume global Lipschitzness of the full strongly convex \(U\).



\texttt{lipschitz\_\allowbreak{}full\_\allowbreak{}potential\_\allowbreak{}not\_\allowbreak{}proximal\_\allowbreak{}deficit}: the source steps apply or rewrite \texttt{bounded\_\allowbreak{}field\_\allowbreak{}not\_\allowbreak{}strongly\allowbreak{}Monotone}, \texttt{Corrections.\allowbreak{}proximal\_\allowbreak{}displacement\_\allowbreak{}strong\_\allowbreak{}monotonicity}.
\end{leanbox}

\end{document}
